\section{Results and Discussion by Research Question}
\label{sec:results}

Results are organized by research question. All reported values come from
frozen experiment artifacts; Online Resource~1, Note~S6, lists the integrity
hashes of those included in the code release.

\subsection{RQ1 (representation and generalization)}
\label{sec:rq1}

RQ1 is answered by the 18 formal runs of Sect.~\ref{sec:setup-parser} on
the standard and character-disjoint splits (Table~\ref{tab:segmentation})
and by the qualitative comparison of
Fig.~\ref{fig:segmentation-qualitative}.

\begin{table*}[t]
\centering
\caption{Segmentation results over three fixed random seeds on the QC-clean
standard split and the character-disjoint split (train, validation, and test
identities disjoint). Values are mean \(\pm\) sample standard deviation;
Endpoint F1 is the strict pixel-level score; bold marks the best value in
each column within a split.}
\label{tab:segmentation}
\footnotesize
\begin{tabular}{llcccc}
\toprule
Split & Model & Macro Dice & Macro IoU & Endpoint F1 & Boundary F1 \\
\midrule
\multirow{3}{*}{Standard}
  & U-Net
  & \(0.9196 \pm 0.0018\)
  & \(0.8515 \pm 0.0030\)
  & \(\mathbf{0.7562 \pm 0.0038}\)
  & \(0.8009 \pm 0.0023\) \\
  & DeepLabV3+
  & \(\mathbf{0.9630 \pm 0.0006}\)
  & \(\mathbf{0.9287 \pm 0.0012}\)
  & \(0.7218 \pm 0.0012\)
  & \(\mathbf{0.8443 \pm 0.0028}\) \\
  & SegFormer-B2
  & \(0.9581 \pm 0.0031\)
  & \(0.9196 \pm 0.0057\)
  & \(0.7139 \pm 0.0046\)
  & \(0.8259 \pm 0.0088\) \\
\midrule
\multirow{3}{*}{Character-disjoint}
  & U-Net
  & \(0.8260 \pm 0.0161\)
  & \(0.7120 \pm 0.0225\)
  & \(\mathbf{0.6972 \pm 0.0045}\)
  & \(0.6430 \pm 0.0136\) \\
  & DeepLabV3+
  & \(0.8664 \pm 0.0094\)
  & \(0.7663 \pm 0.0146\)
  & \(0.6755 \pm 0.0031\)
  & \(0.6728 \pm 0.0138\) \\
  & SegFormer-B2
  & \(\mathbf{0.8866 \pm 0.0008}\)
  & \(\mathbf{0.7977 \pm 0.0013}\)
  & \(0.6604 \pm 0.0033\)
  & \(\mathbf{0.6851 \pm 0.0046}\) \\
\bottomrule
\end{tabular}
\end{table*}


On the standard split, DeepLabV3+ obtains the highest direction Macro Dice
(\(0.9630\pm0.0006\)), Macro IoU (\(0.9287\pm0.0012\)), and Boundary F1
(\(0.8443\pm0.0028\)). U-Net is weaker on direction and boundary
metrics yet gives the best strict endpoint F1 (\(0.7562\pm0.0038\)). Dense
direction segmentation and sparse endpoint localization therefore favour
different model configurations.



All models degrade when character identities are held out
(Table~\ref{tab:segmentation}, bottom). SegFormer-B2
gives the strongest character-disjoint direction result: Macro Dice
\(0.8866\pm0.0008\), Macro IoU \(0.7977\pm0.0013\), and Boundary F1
\(0.6851\pm0.0046\). Its Macro-Dice decrease from the standard split is 7.150
percentage points, versus 9.664 for DeepLabV3+ and 9.360 for U-Net.
DeepLabV3+ remains second for direction overlap, and U-Net retains the
highest strict endpoint F1 (\(0.6972\pm0.0045\)). Near-saturated in-corpus
overlap thus does not imply transfer to unseen character structures, and the
best-performing configuration on held-out identities differs from the
standard-split leader. These comparisons include the initialization and
training choices of Sect.~\ref{sec:setup-parser}, rather than isolating
architecture alone.

\begin{figure*}[!t]
  \centering
  \includegraphics[width=\textwidth]{figures/redrawn/figure4_segmentation_qualitative.pdf}
  \caption{Frozen qualitative comparison for seed 20260811 on a
  multi-channel crossing, an endpoint-rich case, and a held-out character
  (samples are identified as identity/sample index).
  The first two rows use standard-split models and the third uses
  character-disjoint models. Call-outs are identical across columns: the
  circle in (a) marks a crossing where every model emits two labels, the box
  in (b) two short strokes whose endpoints every model keeps, and the boxes
  in (c) black smears and swapped labels on the held-out identity. Values
  under (c) are the dataset-level character-disjoint Macro Dice of
  Table~\ref{tab:segmentation}. Colours follow
  Fig.~\ref{fig:channel-definition}; thresholds are validation-calibrated.}
  \label{fig:segmentation-qualitative}
\end{figure*}

Figure~\ref{fig:segmentation-qualitative} makes the output contract
visible: at the crossing in the first row every model emits overlapping
labels rather than choosing one family, and the third row shows the errors
characteristic of an identity absent from training.

\medskip
\noindent\textbf{Answer to RQ1.}
The best model configuration depends on the target and split:
DeepLabV3+ leads in-corpus direction parsing, SegFormer-B2 unseen-character
direction parsing, and U-Net strict endpoint localization. Holding out
identities costs 7.150--9.664 Macro-Dice points. Crossing agreement is shown
qualitatively, not measured separately; character-disjoint testing measures
transfer within one acquisition source, not to unseen writers.

\subsection{RQ2 (registration and score behaviour)}
\label{sec:rq2}

RQ2 is answered by the controlled-perturbation summary
(Table~\ref{tab:controlled}), the paired alignment ablation
(Table~\ref{tab:alignment}, Fig.~\ref{fig:alignment-ablation}), the inactive-channel audit, and the cross-reference test (Online Resource~1,
Table~S4) described in Sect.~\ref{sec:setup-registration}.

\begin{table*}[t]
\caption{Controlled perturbation summary. Nuisance rows report absolute score
drop; structural rows report signed drop from identity. Full distribution and
severity statistics are in Supplementary Table S1.}
\label{tab:controlled}
\centering
\scriptsize
\begin{tabular}{lrrrrr}
\toprule
Perturbation & Valid \(N\) & Mean drop & 95\% CI & Invalid fraction & Monotonic \\
\midrule
Translation & 800 & 0.000 & [0.000, 0.000] & 0.000 & 1.000 \\
Rotation & 520 & 5.395 & [5.079, 5.747] & 0.133 & 0.647 \\
Scale up & 408 & 13.653 & [12.964, 14.420] & 0.490 & 0.271 \\
Scale down & 800 & 16.356 & [15.670, 17.074] & 0.000 & 0.500 \\
Compound allowed transform & 276 & 8.143 & [7.621, 8.717] & 0.540 & 0.466 \\
Terminal deletion & 800 & 19.757 & [17.387, 22.087] & 0.000 & 0.992 \\
Added direction fragment & 800 & 11.398 & [9.934, 12.929] & 0.000 & 0.987 \\
Local fragment shift & 800 & 4.310 & [3.886, 4.772] & 0.000 & 0.985 \\
Direction dilation & 800 & 24.034 & [22.180, 25.955] & 0.000 & 1.000 \\
Direction erosion & 800 & 22.225 & [20.269, 24.344] & 0.000 & 0.997 \\
Keypoint shift & 800 & 7.468 & [7.294, 7.641] & 0.000 & 1.000 \\
\bottomrule
\end{tabular}
\end{table*}


Under the production registration, translation is recovered exactly in all
800 valid cases, whereas residual sensitivity remains for the other
transforms: mean penalties are 5.395 points for rotation, 13.653 for scale
up, and 16.356 for scale down. Scale up and the compound transform have
invalid fractions of 49.0\% and 54.0\% because of clipping risk, so their
summaries hold only for clipping-safe references and establish no invariance
across the library. Among structural perturbations, direction dilation
and erosion produce the largest penalties, followed by terminal deletion and
added fragments; local fragment shifts and endpoint shifts are weaker but
monotonic in severity.

\begin{table*}[t]
\caption{Paired alignment comparison. Positive benefit differences favour the
current constrained alignment and are computed from unrounded means. Ours is
the production constrained search of Sect.~\ref{sec:alignment-method};
Comparator is no alignment (none) or the wider search (wide); values are
mean score drops; RBC is the rank-biserial correlation. Wilcoxon tests use
per-reference mean paired differences.}
\label{tab:alignment}
\centering
\footnotesize
\begin{tabular}{llrrrrrrr}
\toprule
Comparison & Family & Refs & Ours & Comparator & Benefit 95\% CI & $p$ & RBC \\
\midrule
Ours--none & nuisance & 200 & 8.455 & 38.573 & [29.506, 30.727] & \(<0.001\) & 1.000 \\
Ours--none & structural & 200 & 14.865 & 3.677 & [10.094, 12.322] & \(<0.001\) & 1.000 \\
Ours--wide & nuisance & 200 & 8.455 & 11.351 & [2.757, 3.034] & \(<0.001\) & 1.000 \\
Ours--wide & structural & 200 & 14.865 & 15.149 & [-0.355, -0.213] & \(<0.001\) & -0.683 \\
\bottomrule
\end{tabular}
\end{table*}


The paired comparison isolates what registration buys. Relative to no
alignment, the constrained search reduces the nuisance penalty by 30.118
points (reference-bootstrap 95\% CI 29.506--30.727) and increases the penalty
for the tested structural interventions by 11.189 points (10.094--12.322);
both per-reference signed-rank comparisons are nonzero with rank-biserial
effect size 1.0. Without alignment the same structural edits cost only 3.677 points on
average against 14.865 under registration
(Fig.~\ref{fig:alignment-ablation}c), so the constrained search does not
suppress structural penalties; part of the increase may reflect the centroid-derived translation re-fitting
the whole reference after a local edit. The wider, coarser search incurs a 2.896-point higher nuisance
penalty (2.757--3.034) while preserving only 0.284
additional points of structural penalty. This tested search configuration
does not improve the measured trade-off (perturbation-level results in
Online Resource~1, Table~S2).

\begin{figure*}[!t]
  \centering
  \includegraphics[width=\textwidth]{figures/redrawn/figure5_alignment_ablation.pdf}
  \caption{Alignment ablation on deterministic translation, scale-down, and
  terminal-deletion cases of one reference. Columns compare no alignment, the
  production range, and the wider search; dark grey, blue (dashed outline),
  and red (solid outline) denote overlap, missing reference ink, and extra
  candidate ink. The right-hand column plots the score drop (100 minus
  score) of each condition on one 0--100 axis, with an arrow from no
  alignment to the production range; rows (a)--(b) are nuisance
  perturbations whose drop should vanish, row (c) a structural edit whose
  drop should remain. Registration removes the translation penalty in row (a), reduces but does not eliminate
  the scale penalty in row (b), and retains the structural penalty in
  row (c). Source reference: Calli-Tongji open subset
  \cite{yunmo_jixin_2025}, CC~BY-NC~4.0.}
  \label{fig:alignment-ablation}
\end{figure*}


Forty-one
of the 200 references contain inactive direction channels (36 with one and
five with two; Online Resource~1, Table~S3). Both-empty agreement affects 1,252 of
the 7,604 valid perturbation rows, and averaging over active channels changes
\(S_{\mathrm{prod}}\) by \(-0.289\) points on average and by at most
\(-15.402\): macro overlap can reward absent structure, which motivates the
coverage-aware score tested in RQ3.



In the cross-reference test (Online Resource~1, Table~S4), the seven
same-character cross-style pairs score 20.239 on average (95\% CI 15.639--25.069) against 9.741
(8.584--10.970) for the 50 different-character negatives, with Cliff's
\(\delta=0.84\). The score carries same-character signal across styles, but the small
positive set and the missing same-style condition preclude a broader claim.

\medskip
\noindent\textbf{Answer to RQ2.}
Constrained registration reduces nuisance penalties while retaining
sensitivity to structural edits. The tested wider, coarser search does not
improve this trade-off. Residual rotation and scale penalties and clipping
exclusions bound the claim. The cross-reference test shows same-character
signal, while the inactive-channel audit exposes credit for absent
structure, motivating explicit coverage semantics.

\subsection{RQ3 (human validity)}
\label{sec:rq3}

RQ3 is answered on the 150-pair human-rating cohort of
Sect.~\ref{sec:data-human} with the statistics of
Sect.~\ref{sec:setup-human}: association of the frozen scores
(Table~\ref{tab:human-validation}), rater reliability
(Table~\ref{tab:rater-reliability}), ASDS development
(Table~\ref{tab:spatial-score-development}), and the direct-ink control
(Table~\ref{tab:direct-ink-asds}, Fig.~\ref{fig:asds-direct-ink}).

\begin{table}[t]
\caption{Association between the structural scores and the mean rating of
three blinded human raters across 150 canonical pairs. Confidence intervals
use 10,000 bootstrap resamples clustered by character.}
\label{tab:human-validation}
\centering
\scriptsize
\setlength{\tabcolsep}{3pt}
\begin{tabular}{lccc}
\toprule
Score & \(\rho\) & 95\% CI & \(p\) \\
\midrule
Production & 0.297 & [0.148, 0.441] & \(2.28\times10^{-4}\) \\
Coverage-aware & 0.429 & [0.313, 0.535] & \(4.32\times10^{-8}\) \\
\bottomrule
\end{tabular}
\end{table}

\begin{table}[t]
\caption{Calligraphy-domain-rater reliability. ICC uses only the 150 canonical
presentations. Repeat statistics use the 15 hidden repeats per rater: Exact
and Within \(\pm1\) are agreement fractions between canonical and repeat
ratings, MAD the mean absolute difference, and QWK the quadratic weighted
kappa; E01--E03 are blinded rater codes and POOLED aggregates all repeats.}
\label{tab:rater-reliability}
\centering
\small
\begin{tabular}{lrrrr}
\toprule
Rater & Exact & Within $\pm 1$ & MAD & QWK \\
\midrule
E01 & 0.600 & 1.000 & 0.400 & 0.559 \\
E02 & 0.467 & 1.000 & 0.533 & 0.721 \\
E03 & 0.400 & 0.800 & 0.800 & 0.591 \\
POOLED & 0.489 & 0.933 & 0.578 & 0.634 \\
\bottomrule
\end{tabular}
\vspace{2pt}
\begin{minipage}{0.96\columnwidth}\footnotesize
Canonical ICC(2,1)=0.448 and
ICC(2,$k$)=0.709 for
150 pairs and 3 raters.
\end{minipage}
\end{table}


The frozen production score correlates modestly with the three-rater mean
(\(\rho=0.297\), character-cluster 95\% CI 0.148--0.441), and the
coverage-aware score is stronger (\(\rho=0.429\), 0.313--0.535); their
paired difference is 0.132 (0.062--0.197). Single-rater reliability is below the 0.50 threshold that Koo and Li
\cite{koo2016icc} label moderate (ICC(2,1) 0.448), whereas the three-rater
mean is moderate (ICC(2,\(k\)) 0.709); hidden repeats yield
48.9\% exact agreement, 93.3\% agreement within one point, MAD 0.578, and
pooled QWK 0.634.

\begin{table}[t]
\caption{Development-stage ASDS association with mean human rating
(character-cluster bootstrap 95\% CI).}
\label{tab:spatial-score-development}
\centering
\scriptsize
\setlength{\tabcolsep}{3pt}
\begin{tabular}{@{}lcc@{}}
\toprule
Score or component & Spearman \(\rho\) & 95\% CI \\
\midrule
Grid only & 0.447 & [0.303, 0.572] \\
Projection only & 0.445 & [0.316, 0.564] \\
Polar only & 0.547 & [0.443, 0.645] \\
ASDS (frozen) & 0.556 & [0.454, 0.651] \\
ASDS (grouped OOF) & 0.540 & [0.432, 0.639] \\
\bottomrule
\end{tabular}
\end{table}


ASDS reaches \(\rho=0.556\) (0.454--0.651), and character-grouped
out-of-fold weight selection retains \(\rho=0.540\) (0.432--0.639). The
paired intervals for ASDS minus the production and coverage-aware scores are
0.110--0.410 and 0.017--0.239; the out-of-fold difference from the
coverage-aware score is less certain (\(-0.003\)--0.224). The polar component
is the strongest single term (\(\rho=0.547\)); grid and projection reach
0.447 and 0.445. Because the same ratings informed feature design, \(\rho=0.556\) is a
retrospective development result, and \(\rho=0.540\) is a
character-grouped internal-validation estimate that removes
weight-selection optimism only: the feature families were also chosen after
the ratings, so neither value is an external or prospective confirmation.

\begin{table}[t]
\caption{Direct-ink ablation of the frozen ASDS specification on the same 150
development pairs.}
\label{tab:direct-ink-asds}
\centering
\begin{tabular}{@{}lcc@{}}
\toprule
ASDS input & Spearman $\rho$ & 95\% CI \\
\midrule
Annotated-mask union & 0.556 & [0.454, 0.651] \\
Direct ink & 0.556 & [0.454, 0.649] \\
\midrule
$\Delta\rho$ (annotated--direct) & 0.000 & [-0.010, 0.010] \\
\bottomrule
\end{tabular}
\parbox{0.86\columnwidth}{\footnotesize\raggedright
Direct ink: grayscale intensity $<240$ (frozen QC threshold,
Sect.~\ref{sec:setup-human}); same registration, features, and weights for
both variants.}
\end{table}


The direct-ink control tests whether an annotated-mask union improves ASDS
association over thresholded source ink. Annotated-mask-union and direct-ink ASDS reach
\(\rho=0.55623\) and \(0.55583\);
their difference is \(\Delta\rho=0.00040\) (95\% CI
\(-0.01011\)--\(0.01030\); Table~\ref{tab:direct-ink-asds} rounds to
three decimals), and the two score vectors have rank association 0.9973
(Fig.~\ref{fig:asds-direct-ink}). On these clean direct-digital development pairs, thresholded ink and the
annotated-mask union yield closely matched human associations and score
rankings. This result supports the silhouette interpretation of ASDS, but
does not establish formal equivalence or transfer to degraded input.

\begin{figure*}[p]
  \centering
  \includegraphics[width=\textwidth]{figures/revision/figure6_asds_direct_ink.pdf}
  \caption{ASDS mechanism and direct-ink control. (a)--(b)~A natural pair
  and its aligned silhouettes. (c)--(d)~Absolute polar (four radial by eight
  angular bins) and \(3\times3\) occupancy differences, each normalized
  to its own maximum; intensities are not comparable between maps.
  (e)~Projection profiles: rows solid, columns dashed; candidate red,
  reference blue. (f)~Frozen formula and example score. (g)~All 150 pairs
  from 40 identities. Annotated-mask union uses five manual direction masks,
  not parser predictions. Brackets are character-cluster bootstrap 95\% CIs;
  \(\Delta\rho\) is annotated-mask-union minus direct-ink association.
  Filled points identify the example; the right-hand dashed line is the
  identity. Observed associations and rankings are closely matched.}
  \label{fig:asds-direct-ink}
\end{figure*}

\medskip
\noindent\textbf{Answer to RQ3.}
ASDS has the strongest observed human association, but its internally
validated advantage over the coverage-aware score remains unresolved
(paired 95\% CI \(-0.003\)--0.224). Its silhouette formulation and the
direct-ink control separate scalar agreement from the parser's role in
direction and endpoint evidence. RQ1 evaluates parser accuracy; RQ4 tests
the downstream diagnostic mechanism under controlled mask interventions.
Human validity remains limited to structural agreement on one retrospective
cohort with three raters, two of them authors.

\subsection{RQ4 (localized diagnosis)}
\label{sec:rq4}

RQ4 is answered on the diagnostic cohort with the metrics of
Sect.~\ref{sec:setup-diagnostic} (Table~\ref{tab:feedback-failures},
Fig.~\ref{fig:diagnostic-cases}).

\begin{table}[t]
\caption{Primary taxonomy of exact-region diagnostic failures.}
\label{tab:feedback-failures}
\centering
\small
\begin{tabular}{lrr}
\toprule
Failure type & Count & Fraction \\
\midrule
multi\_region\_or\_grid\_boundary\_truth & 1922 & 0.539 \\
wrong\_direction\_channel & 876 & 0.246 \\
local\_finding\_absent\_from\_top3 & 718 & 0.201 \\
wrong\_missing\_extra\_type & 38 & 0.011 \\
alignment\_residual\_region\_shift & 11 & 0.003 \\
\bottomrule
\end{tabular}
\end{table}


Required Recall@3 is 0.701, whereas exact grid localization is 0.109. Under the fixed-order taxonomy, 53.9\% of exact-region failures are assigned
to multi-cell or grid-boundary truth, 24.6\% to wrong direction attribution,
and 20.1\% to a local finding absent from the top three. Figure~\ref{fig:diagnostic-cases} shows one exact
success and one audited failure of the first kind: the returned cell
overlaps the multi-cell change, and the frozen metric still records a
failure. The largest assigned category therefore reflects a mismatch between
spatial extent and the exact-cell metric, followed by direction attribution.
These categories concern structured evidence rather than wording.

\begin{figure*}[p]
  \centering
  \includegraphics[width=\textwidth]{figures/redrawn/figure7_diagnostic_cases.pdf}
  \caption{Frozen diagnostic evidence: (a)~an exact-localization success
  and (b)~multi-cell truth classified as failure by the exact-one-cell metric.
  Findings report type~$|$ direction~$|$ missing/extra~$|$ cell; rXcY indexes
  the \(3\times3\) grid and \texttt{user} denotes the candidate.
  Dark grey is overlap, blue (dashed) missing reference ink, and red (solid)
  extra candidate ink. Solid and dashed black frames mark truth and predicted
  cells, respectively. Cards state the findings, truth, and audit verdict.
  The bottom strip summarizes Table~\ref{tab:feedback-failures} over
  3,565 exact-region failures; case (b) is in the largest assigned category.}
  \label{fig:diagnostic-cases}
\end{figure*}

\medskip
\noindent\textbf{Answer to RQ4.}
The evidence layer usually surfaces the required finding but does not
reliably assign it to one \(3\times3\) cell. Multi-cell truth and direction
attribution dominate the assigned failures. Wrong missing/extra type is the
primary category in 1.1\%, a lower bound under the fixed-order taxonomy.
These controlled interventions evaluate the diagnostic mechanism, not
teacher agreement or pedagogical usefulness.
